\begin{center}
    \section*{\fontsize{20}{20}\selectfont Appendix A}
\end{center}
\vspace{10mm}

\section*{Engineering Standards Compliance}

This appendix provides detailed documentation of how the research addresses complex engineering problems and activities as defined by accreditation standards (ABET, BAETE). This analysis demonstrates the depth, interdisciplinary integration, and societal impact of the work.

\subsection*{Complex Engineering Problems (P1-P7)}

\subsubsection*{P1: Depth of Knowledge Required}

This research requires advanced understanding across multiple specialized domains:

\textbf{Graph Neural Networks:} Message-passing architectures (GraphSAGE, GAT), inductive learning, neighborhood sampling strategies, and graph attention mechanisms.

\textbf{Medical Image Processing:} SLIC superpixel generation, 3D MRI preprocessing (skull-stripping, normalization, registration), multi-modal intensity analysis.

\textbf{Deep Learning Optimization:} Adam optimizer, batch normalization, gradient accumulation, learning rate scheduling, early stopping strategies.

\textbf{Software Engineering:} PyTorch Geometric data structures, CUDA optimization, medical imaging standards (NIfTI, DICOM), version control and reproducibility protocols.

\subsubsection*{P2: Range of Conflicting Requirements}

The project navigates multiple competing objectives:

\begin{itemize}
    \item \textbf{Accuracy vs. Speed:} Target $>$90\% Dice coefficient while maintaining $<$50ms inference time
    \item \textbf{Expressiveness vs. Efficiency:} Sufficient model capacity for tumor detection within $<$1M parameter budget
    \item \textbf{Feature Richness vs. Construction Speed:} 15 engineered features capturing tumor characteristics without excessive preprocessing overhead
    \item \textbf{Validation Rigor vs. Iteration Speed:} Comprehensive 5-fold cross-validation with integrity checks while maintaining reasonable experimentation cycle time
\end{itemize}

\subsubsection*{P3: Depth of Analysis Required}

Extensive experimental validation conducted:

\begin{itemize}
    \item Stratified 5-fold cross-validation with patient-level splits (1,000 train, 251 test per fold)
    \item Systematic ablation studies: 5 vs. 6 layers, batch sizes 32/48/64
    \item Hyperparameter tuning across 20+ configurations (learning rate, hidden dimensions, dropout)
    \item Comprehensive integrity auditing: 15 validation checks ensuring zero data leakage
    \item Multi-metric evaluation: Dice, sensitivity, specificity, precision, inference time, memory, parameters
\end{itemize}

\subsubsection*{P4: Familiarity with Domain-Specific Issues}

Critical medical imaging challenges addressed:

\begin{itemize}
    \item \textbf{Protocol Variability:} 1,251 multi-institutional MRI scans with varying acquisition parameters
    \item \textbf{Image Artifacts:} Noise, motion artifacts, intensity inhomogeneities requiring robust preprocessing
    \item \textbf{Class Imbalance:} Extreme imbalance (tumors occupy 5-10\% of brain volume) requiring weighted loss functions
    \item \textbf{Privacy \& Ethics:} Patient data handling following HIPAA/GDPR regulations
    \item \textbf{Clinical Consequences:} False negatives can delay treatment; false positives cause unnecessary interventions
    \item \textbf{Data Leakage:} Zero ground-truth leakage critical for valid medical AI (removed tumor\_ratio feature)
\end{itemize}

\subsubsection*{P5: Extent of Applicable Standards and Codes}

Compliance with multiple standards:

\begin{itemize}
    \item \textbf{Medical Imaging Formats:} NIfTI (research), DICOM (clinical deployment)
    \item \textbf{Data Privacy:} HIPAA (USA), GDPR (Europe) for patient data protection
    \item \textbf{Evaluation Protocols:} BraTS challenge metrics and evaluation procedures
    \item \textbf{Reproducibility Standards:} Seed management, deterministic CUDA operations, dependency versioning
    \item \textbf{Custom Validation:} Novel integrity auditing procedures where established benchmarks for GNN medical segmentation did not exist
\end{itemize}

\subsubsection*{P6: Level of Uncertainty}

High unpredictability from multiple sources:

\begin{itemize}
    \item \textbf{Tumor Heterogeneity:} Extreme variation in glioma appearance (shape, size, intensity, texture)
    \item \textbf{Scanner Variability:} Different magnetic field strengths (1.5T, 3T), acquisition protocols, vendor hardware (GE, Siemens, Philips)
    \item \textbf{Patient Factors:} Motion during scanning, anatomical variations, pathological diversity
    \item \textbf{Architectural Uncertainty:} Optimal graph construction strategy unclear (superpixel size, feature engineering choices)
    \item \textbf{Generalization Risk:} Training on 1,251 cases may not cover full diversity of clinical presentations
\end{itemize}

\subsubsection*{P7: Component Interdependence}

Tightly coupled system requiring orchestration:

\begin{itemize}
    \item \textbf{Preprocessing $\rightarrow$ Features:} Skull-stripping and normalization directly impact feature quality
    \item \textbf{Graph Construction $\rightarrow$ Learning:} SLIC superpixel segmentation and adjacency determination define model input structure
    \item \textbf{Architecture $\rightarrow$ Performance:} Network depth, hidden dimensions depend on graph properties
    \item \textbf{Training Protocol $\rightarrow$ Convergence:} Batch size, learning rate, augmentation affect final accuracy
    \item \textbf{Ensemble Strategy $\rightarrow$ Diversity:} Individual model variance must be sufficient for ensemble gains
    \item \textbf{Clinical Integration:} Potential connection to visualization tools (3D Slicer), PACS systems requires standardized outputs
\end{itemize}

\subsection*{Complex Engineering Activities (A1-A5)}

\subsubsection*{A1: Diversity of Resources}

Integration of diverse computational and domain resources:

\begin{itemize}
    \item \textbf{Dataset:} BraTS 2021 (1,251 multi-parametric MRI volumes, expert annotations)
    \item \textbf{Software Libraries:} PyTorch Geometric (GNN), SimpleITK (medical imaging), scikit-image (superpixels)
    \item \textbf{Hardware:} NVIDIA RTX 2060 6GB VRAM (consumer-grade, demonstrating accessibility)
    \item \textbf{Domain Expertise:} Neuroimaging literature (glioma biology, MRI physics)
    \item \textbf{Validation Frameworks:} ML best practices (cross-validation, ablation, ensembles)
\end{itemize}

\subsubsection*{A2: Level of Interdisciplinary Interaction}

Continuous collaboration across fields:

\begin{itemize}
    \item Understanding clinical requirements (radiologist consultation for accuracy/speed needs)
    \item Translating medical knowledge into graph features and network design
    \item Ensuring data integrity through validation protocols preventing medical AI pitfalls
    \item Interpreting predictions in clinical context (sensitivity to false negatives vs. positives)
    \item Balancing computer science practices (reproducibility) with medical standards (privacy, validation)
    \item Communicating efficiency-accuracy trade-offs to diverse stakeholders
\end{itemize}

\subsubsection*{A3: Innovation and Novelty}

Novel technical contributions:

\begin{itemize}
    \item \textbf{Superpixel-GNN Integration:} First optimized graph construction for 3D medical images with 15 multi-modal features
    \item \textbf{Efficiency Study:} Comprehensive GNN efficiency vs. accuracy trade-off analysis for brain segmentation
    \item \textbf{Batch Size Discovery:} Characterization of batch size sensitivity in medical imaging GNN tasks (32 optimal, 64 degrades)
    \item \textbf{Architectural Validation:} Systematic ablation showing 5-layer sweet spot (6 layers no benefit)
    \item \textbf{Integrity Framework:} 15 comprehensive checks addressing reproducibility crisis in medical AI
    \item \textbf{Performance Demonstration:} CNN-competitive accuracy (92.92\%) with 6.9$\times$ speedup, 156$\times$ parameter reduction
\end{itemize}

\subsubsection*{A4: Societal and Environmental Consequences}

Significant implications for healthcare and sustainability:

\begin{itemize}
    \item \textbf{Healthcare Access:} Enables deployment in developing countries, rural clinics, mobile diagnostic units where volumetric CNNs impractical
    \item \textbf{Clinical Efficiency:} 6.9$\times$ faster inference reduces patient waiting, enables real-time intra-operative guidance
    \item \textbf{Energy Sustainability:} 156$\times$ smaller models reduce carbon footprint at scale
    \item \textbf{Clinical Applications:} Improved segmentation assists diagnosis, treatment planning (surgery, radiation), therapeutic monitoring
    \item \textbf{Error Consequences:} Direct clinical impact requiring careful validation and uncertainty quantification
    \item \textbf{Privacy Protection:} Ethical patient data handling through privacy-preserving protocols
    \item \textbf{Global Impact:} Contribution to democratizing AI-assisted medical diagnosis worldwide
\end{itemize}

\subsubsection*{A5: Knowledge Application and Self-Directed Learning}

Autonomous learning beyond standard curriculum:

\begin{itemize}
    \item \textbf{Medical Imaging:} MRI physics, T1/T2/FLAIR contrasts, clinical interpretation
    \item \textbf{Neuroanatomy:} Glioma pathology, tumor growth patterns, anatomical landmarks
    \item \textbf{Medical Formats:} NIfTI structures, skull-stripping algorithms, intensity normalization
    \item \textbf{Graph Theory:} Message passing, graph convolutions, attention mechanisms
    \item \textbf{Medical AI Validation:} Patient-level splitting, integrity checks, clinical metrics
    \item \textbf{HPC Optimization:} GPU memory management, mixed-precision training, distributed computing
    \item \textbf{Communication:} Bridging computer science, medical imaging, and clinical practice
\end{itemize}

\vspace{5mm}

This appendix demonstrates that the research addresses complex engineering challenges requiring depth across multiple disciplines, innovation in methodology, and consideration of societal impacts. The work exemplifies engineering problem-solving at the intersection of computer science, medical imaging, and healthcare delivery.

\clearpage
